% Sinh một lần từ cấu trúc notebook; mỗi mục có đủ 6 tiểu mục theo đặc tả (Bài 11).

\subsection{11.1\quad Chuỗi thời gian, dừng và nhiễu trắng}
\textit{Bộ dữ liệu: HU, LY, DR (cả 3)}

\subsubsection*{Nội dung lý thuyết cần kiểm chứng}
\TODO{viết ngắn gọn, chính xác theo bài giảng, có số hiệu Định nghĩa/Định lý/Mệnh đề (xem ô ``Nội dung lý thuyết'' của notebook).}

\subsubsection*{Cách triển khai}
\TODO{áp dụng trên bộ nào, kiểu [D] hay [M] (xem ô ``Cách triển khai'').}

\subsubsection*{Mã lệnh Python}
\TODO{dán đoạn mã quan trọng dạng văn bản trong khung \texttt{lstlisting} (không chụp ảnh).}

\subsubsection*{Kết quả tính toán}
\TODO{chèn bảng/hình từ \texttt{reports/tables} và \texttt{reports/figures} do notebook sinh ra.}

\subsubsection*{Phân tích và diễn giải}
\TODO{giải thích ý nghĩa thống kê và liên hệ trực tiếp với lý thuyết.}

\subsubsection*{Nhận xét}
\TODO{phù hợp lý thuyết hay không; nếu lệch thì nguyên nhân hợp lý.}

\subsection{11.2\quad Tính chất của hàm tự hiệp phương sai}
\textit{Bộ dữ liệu: HU, LY, DR (cả 3)}

\subsubsection*{Nội dung lý thuyết cần kiểm chứng}
\TODO{viết ngắn gọn, chính xác theo bài giảng, có số hiệu Định nghĩa/Định lý/Mệnh đề (xem ô ``Nội dung lý thuyết'' của notebook).}

\subsubsection*{Cách triển khai}
\TODO{áp dụng trên bộ nào, kiểu [D] hay [M] (xem ô ``Cách triển khai'').}

\subsubsection*{Mã lệnh Python}
\TODO{dán đoạn mã quan trọng dạng văn bản trong khung \texttt{lstlisting} (không chụp ảnh).}

\subsubsection*{Kết quả tính toán}
\TODO{chèn bảng/hình từ \texttt{reports/tables} và \texttt{reports/figures} do notebook sinh ra.}

\subsubsection*{Phân tích và diễn giải}
\TODO{giải thích ý nghĩa thống kê và liên hệ trực tiếp với lý thuyết.}

\subsubsection*{Nhận xét}
\TODO{phù hợp lý thuyết hay không; nếu lệch thì nguyên nhân hợp lý.}

\subsection{11.3\quad Trung bình mẫu khi có phụ thuộc}
\textit{Bộ dữ liệu: HU (chính), LY, DR; mô phỏng AR(1)}

\subsubsection*{Nội dung lý thuyết cần kiểm chứng}
\TODO{viết ngắn gọn, chính xác theo bài giảng, có số hiệu Định nghĩa/Định lý/Mệnh đề (xem ô ``Nội dung lý thuyết'' của notebook).}

\subsubsection*{Cách triển khai}
\TODO{áp dụng trên bộ nào, kiểu [D] hay [M] (xem ô ``Cách triển khai'').}

\subsubsection*{Mã lệnh Python}
\TODO{dán đoạn mã quan trọng dạng văn bản trong khung \texttt{lstlisting} (không chụp ảnh).}

\subsubsection*{Kết quả tính toán}
\TODO{chèn bảng/hình từ \texttt{reports/tables} và \texttt{reports/figures} do notebook sinh ra.}

\subsubsection*{Phân tích và diễn giải}
\TODO{giải thích ý nghĩa thống kê và liên hệ trực tiếp với lý thuyết.}

\subsubsection*{Nhận xét}
\TODO{phù hợp lý thuyết hay không; nếu lệch thì nguyên nhân hợp lý.}

\subsection{11.4\quad ACF mẫu, dải iid, Bartlett và kiểm định Ljung–Box}
\textit{Bộ dữ liệu: HU, LY, DR (cả 3) + mô phỏng nhiễu iid}

\subsubsection*{Nội dung lý thuyết cần kiểm chứng}
\TODO{viết ngắn gọn, chính xác theo bài giảng, có số hiệu Định nghĩa/Định lý/Mệnh đề (xem ô ``Nội dung lý thuyết'' của notebook).}

\subsubsection*{Cách triển khai}
\TODO{áp dụng trên bộ nào, kiểu [D] hay [M] (xem ô ``Cách triển khai'').}

\subsubsection*{Mã lệnh Python}
\TODO{dán đoạn mã quan trọng dạng văn bản trong khung \texttt{lstlisting} (không chụp ảnh).}

\subsubsection*{Kết quả tính toán}
\TODO{chèn bảng/hình từ \texttt{reports/tables} và \texttt{reports/figures} do notebook sinh ra.}

\subsubsection*{Phân tích và diễn giải}
\TODO{giải thích ý nghĩa thống kê và liên hệ trực tiếp với lý thuyết.}

\subsubsection*{Nhận xét}
\TODO{phù hợp lý thuyết hay không; nếu lệch thì nguyên nhân hợp lý.}

\subsection{11.5\quad Quá trình tuyến tính, toán tử lùi, trọng số $\psi$ và ACF lý thuyết}
\textit{Bộ dữ liệu: HU, LY (DR: giải thích không áp dụng)}

\subsubsection*{Nội dung lý thuyết cần kiểm chứng}
\TODO{viết ngắn gọn, chính xác theo bài giảng, có số hiệu Định nghĩa/Định lý/Mệnh đề (xem ô ``Nội dung lý thuyết'' của notebook).}

\subsubsection*{Cách triển khai}
\TODO{áp dụng trên bộ nào, kiểu [D] hay [M] (xem ô ``Cách triển khai'').}

\subsubsection*{Mã lệnh Python}
\TODO{dán đoạn mã quan trọng dạng văn bản trong khung \texttt{lstlisting} (không chụp ảnh).}

\subsubsection*{Kết quả tính toán}
\TODO{chèn bảng/hình từ \texttt{reports/tables} và \texttt{reports/figures} do notebook sinh ra.}

\subsubsection*{Phân tích và diễn giải}
\TODO{giải thích ý nghĩa thống kê và liên hệ trực tiếp với lý thuyết.}

\subsubsection*{Nhận xét}
\TODO{phù hợp lý thuyết hay không; nếu lệch thì nguyên nhân hợp lý.}

\subsection{11.6\quad Mô hình AR(1)}
\textit{Bộ dữ liệu: HU (chính); LY, DR: giải thích}

\subsubsection*{Nội dung lý thuyết cần kiểm chứng}
\TODO{viết ngắn gọn, chính xác theo bài giảng, có số hiệu Định nghĩa/Định lý/Mệnh đề (xem ô ``Nội dung lý thuyết'' của notebook).}

\subsubsection*{Cách triển khai}
\TODO{áp dụng trên bộ nào, kiểu [D] hay [M] (xem ô ``Cách triển khai'').}

\subsubsection*{Mã lệnh Python}
\TODO{dán đoạn mã quan trọng dạng văn bản trong khung \texttt{lstlisting} (không chụp ảnh).}

\subsubsection*{Kết quả tính toán}
\TODO{chèn bảng/hình từ \texttt{reports/tables} và \texttt{reports/figures} do notebook sinh ra.}

\subsubsection*{Phân tích và diễn giải}
\TODO{giải thích ý nghĩa thống kê và liên hệ trực tiếp với lý thuyết.}

\subsubsection*{Nhận xét}
\TODO{phù hợp lý thuyết hay không; nếu lệch thì nguyên nhân hợp lý.}

\subsection{11.7\quad AR(p): tính nhân quả và phương trình Yule–Walker}
\textit{Bộ dữ liệu: LY, HU (DR: không áp dụng)}

\subsubsection*{Nội dung lý thuyết cần kiểm chứng}
\TODO{viết ngắn gọn, chính xác theo bài giảng, có số hiệu Định nghĩa/Định lý/Mệnh đề (xem ô ``Nội dung lý thuyết'' của notebook).}

\subsubsection*{Cách triển khai}
\TODO{áp dụng trên bộ nào, kiểu [D] hay [M] (xem ô ``Cách triển khai'').}

\subsubsection*{Mã lệnh Python}
\TODO{dán đoạn mã quan trọng dạng văn bản trong khung \texttt{lstlisting} (không chụp ảnh).}

\subsubsection*{Kết quả tính toán}
\TODO{chèn bảng/hình từ \texttt{reports/tables} và \texttt{reports/figures} do notebook sinh ra.}

\subsubsection*{Phân tích và diễn giải}
\TODO{giải thích ý nghĩa thống kê và liên hệ trực tiếp với lý thuyết.}

\subsubsection*{Nhận xét}
\TODO{phù hợp lý thuyết hay không; nếu lệch thì nguyên nhân hợp lý.}

\subsection{11.8\quad AR(2): tam giác nhân quả và chu kỳ}
\textit{Bộ dữ liệu: LY (chính), HU; DR: không áp dụng}

\subsubsection*{Nội dung lý thuyết cần kiểm chứng}
\TODO{viết ngắn gọn, chính xác theo bài giảng, có số hiệu Định nghĩa/Định lý/Mệnh đề (xem ô ``Nội dung lý thuyết'' của notebook).}

\subsubsection*{Cách triển khai}
\TODO{áp dụng trên bộ nào, kiểu [D] hay [M] (xem ô ``Cách triển khai'').}

\subsubsection*{Mã lệnh Python}
\TODO{dán đoạn mã quan trọng dạng văn bản trong khung \texttt{lstlisting} (không chụp ảnh).}

\subsubsection*{Kết quả tính toán}
\TODO{chèn bảng/hình từ \texttt{reports/tables} và \texttt{reports/figures} do notebook sinh ra.}

\subsubsection*{Phân tích và diễn giải}
\TODO{giải thích ý nghĩa thống kê và liên hệ trực tiếp với lý thuyết.}

\subsubsection*{Nhận xét}
\TODO{phù hợp lý thuyết hay không; nếu lệch thì nguyên nhân hợp lý.}

\subsection{11.9\quad Mô hình MA(q): ACF cắt cụt, MA(1), khả nghịch}
\textit{Bộ dữ liệu: DR, HU (LY: giải thích) + mô phỏng MA(2)}

\subsubsection*{Nội dung lý thuyết cần kiểm chứng}
\TODO{viết ngắn gọn, chính xác theo bài giảng, có số hiệu Định nghĩa/Định lý/Mệnh đề (xem ô ``Nội dung lý thuyết'' của notebook).}

\subsubsection*{Cách triển khai}
\TODO{áp dụng trên bộ nào, kiểu [D] hay [M] (xem ô ``Cách triển khai'').}

\subsubsection*{Mã lệnh Python}
\TODO{dán đoạn mã quan trọng dạng văn bản trong khung \texttt{lstlisting} (không chụp ảnh).}

\subsubsection*{Kết quả tính toán}
\TODO{chèn bảng/hình từ \texttt{reports/tables} và \texttt{reports/figures} do notebook sinh ra.}

\subsubsection*{Phân tích và diễn giải}
\TODO{giải thích ý nghĩa thống kê và liên hệ trực tiếp với lý thuyết.}

\subsubsection*{Nhận xét}
\TODO{phù hợp lý thuyết hay không; nếu lệch thì nguyên nhân hợp lý.}

\subsection{11.10\quad ARMA(p, q): nhân quả, khả nghịch, nghiệm chung}
\textit{Bộ dữ liệu: HU, LY (DR: minh họa thừa tham số)}

\subsubsection*{Nội dung lý thuyết cần kiểm chứng}
\TODO{viết ngắn gọn, chính xác theo bài giảng, có số hiệu Định nghĩa/Định lý/Mệnh đề (xem ô ``Nội dung lý thuyết'' của notebook).}

\subsubsection*{Cách triển khai}
\TODO{áp dụng trên bộ nào, kiểu [D] hay [M] (xem ô ``Cách triển khai'').}

\subsubsection*{Mã lệnh Python}
\TODO{dán đoạn mã quan trọng dạng văn bản trong khung \texttt{lstlisting} (không chụp ảnh).}

\subsubsection*{Kết quả tính toán}
\TODO{chèn bảng/hình từ \texttt{reports/tables} và \texttt{reports/figures} do notebook sinh ra.}

\subsubsection*{Phân tích và diễn giải}
\TODO{giải thích ý nghĩa thống kê và liên hệ trực tiếp với lý thuyết.}

\subsubsection*{Nhận xét}
\TODO{phù hợp lý thuyết hay không; nếu lệch thì nguyên nhân hợp lý.}

\subsection{11.11\quad Tự tương quan riêng phần (PACF)}
\textit{Bộ dữ liệu: HU (đã khử xu thế), LY, DR (cả 3)}

\subsubsection*{Nội dung lý thuyết cần kiểm chứng}
\TODO{viết ngắn gọn, chính xác theo bài giảng, có số hiệu Định nghĩa/Định lý/Mệnh đề (xem ô ``Nội dung lý thuyết'' của notebook).}

\subsubsection*{Cách triển khai}
\TODO{áp dụng trên bộ nào, kiểu [D] hay [M] (xem ô ``Cách triển khai'').}

\subsubsection*{Mã lệnh Python}
\TODO{dán đoạn mã quan trọng dạng văn bản trong khung \texttt{lstlisting} (không chụp ảnh).}

\subsubsection*{Kết quả tính toán}
\TODO{chèn bảng/hình từ \texttt{reports/tables} và \texttt{reports/figures} do notebook sinh ra.}

\subsubsection*{Phân tích và diễn giải}
\TODO{giải thích ý nghĩa thống kê và liên hệ trực tiếp với lý thuyết.}

\subsubsection*{Nhận xét}
\TODO{phù hợp lý thuyết hay không; nếu lệch thì nguyên nhân hợp lý.}

\subsection{11.12\quad Ước lượng: Yule–Walker, OLS, cực đại hợp lý}
\textit{Bộ dữ liệu: LY, HU (+ mô phỏng AR(2)); DR: không áp dụng}

\subsubsection*{Nội dung lý thuyết cần kiểm chứng}
\TODO{viết ngắn gọn, chính xác theo bài giảng, có số hiệu Định nghĩa/Định lý/Mệnh đề (xem ô ``Nội dung lý thuyết'' của notebook).}

\subsubsection*{Cách triển khai}
\TODO{áp dụng trên bộ nào, kiểu [D] hay [M] (xem ô ``Cách triển khai'').}

\subsubsection*{Mã lệnh Python}
\TODO{dán đoạn mã quan trọng dạng văn bản trong khung \texttt{lstlisting} (không chụp ảnh).}

\subsubsection*{Kết quả tính toán}
\TODO{chèn bảng/hình từ \texttt{reports/tables} và \texttt{reports/figures} do notebook sinh ra.}

\subsubsection*{Phân tích và diễn giải}
\TODO{giải thích ý nghĩa thống kê và liên hệ trực tiếp với lý thuyết.}

\subsubsection*{Nhận xét}
\TODO{phù hợp lý thuyết hay không; nếu lệch thì nguyên nhân hợp lý.}

\subsection{11.13\quad Chọn bậc bằng AIC/BIC và thẩm định cuốn chiếu}
\textit{Bộ dữ liệu: LY, HU (DR: lưới ARMA nhỏ)}

\subsubsection*{Nội dung lý thuyết cần kiểm chứng}
\TODO{viết ngắn gọn, chính xác theo bài giảng, có số hiệu Định nghĩa/Định lý/Mệnh đề (xem ô ``Nội dung lý thuyết'' của notebook).}

\subsubsection*{Cách triển khai}
\TODO{áp dụng trên bộ nào, kiểu [D] hay [M] (xem ô ``Cách triển khai'').}

\subsubsection*{Mã lệnh Python}
\TODO{dán đoạn mã quan trọng dạng văn bản trong khung \texttt{lstlisting} (không chụp ảnh).}

\subsubsection*{Kết quả tính toán}
\TODO{chèn bảng/hình từ \texttt{reports/tables} và \texttt{reports/figures} do notebook sinh ra.}

\subsubsection*{Phân tích và diễn giải}
\TODO{giải thích ý nghĩa thống kê và liên hệ trực tiếp với lý thuyết.}

\subsubsection*{Nhận xét}
\TODO{phù hợp lý thuyết hay không; nếu lệch thì nguyên nhân hợp lý.}

\subsection{11.14\quad Chẩn đoán phần dư và checklist báo cáo ARMA}
\textit{Bộ dữ liệu: HU, LY, DR (cả 3)}

\subsubsection*{Nội dung lý thuyết cần kiểm chứng}
\TODO{viết ngắn gọn, chính xác theo bài giảng, có số hiệu Định nghĩa/Định lý/Mệnh đề (xem ô ``Nội dung lý thuyết'' của notebook).}

\subsubsection*{Cách triển khai}
\TODO{áp dụng trên bộ nào, kiểu [D] hay [M] (xem ô ``Cách triển khai'').}

\subsubsection*{Mã lệnh Python}
\TODO{dán đoạn mã quan trọng dạng văn bản trong khung \texttt{lstlisting} (không chụp ảnh).}

\subsubsection*{Kết quả tính toán}
\TODO{chèn bảng/hình từ \texttt{reports/tables} và \texttt{reports/figures} do notebook sinh ra.}

\subsubsection*{Phân tích và diễn giải}
\TODO{giải thích ý nghĩa thống kê và liên hệ trực tiếp với lý thuyết.}

\subsubsection*{Nhận xét}
\TODO{phù hợp lý thuyết hay không; nếu lệch thì nguyên nhân hợp lý.}
